
\medskip
\noindent\underline{Corollary \SGARefTarget{sga7.sga7i.eviii.n3-5}{3.5}.}  \underline{Let $P$ be a smooth
group scheme of finite presentation over $S$, with connected fibres,
and let}
\[
 0\longrightarrow Q'\longrightarrow Q\longrightarrow Q''
 \longrightarrow0
\]
\underline{be an exact sequence of commutative groups over $S$, with
$Q'$ a torus.  Then the inverse-image functor on biextensions}
\begin{equation}\tag{\SGARefTarget{sga7.sga7i.eviii.n3-5-1}{3.5.1}}
 \operatorname{BIEXT}(P,Q'';\underline G_m)
 \xrightarrow{\ \sim\ }
 \operatorname{BIEXT}(P,Q;\underline G_m)
\end{equation}
\underline{is an equivalence of categories.  In particular, the
following natural homomorphisms are isomorphisms:}
\begin{equation}\tag{\SGARefTarget{sga7.sga7i.eviii.n3-5-2}{3.5.2}}
 \Biext^i(P,Q'';\underline G_m)
 \xrightarrow{\ \sim\ }
 \Biext^i(P,Q;\underline G_m),
 \qquad i=0,1.
\end{equation}

Indeed, up to equivalence, $\operatorname{BIEXT}(P,Q'';\underline
G_m)$ identifies with the category of biextensions $E$ of $(P,Q)$ by
$\underline G_m$ \underline{equipped} with a trivialization of the
induced biextension $E'$ of $(P,Q')$ by $\underline G_m$ (VII
\SGARefLink{sga7.sga7i.evii.n3-7-6}{3.7.6}).  By \SGARefLink{sga7.sga7i.eviii.n3-4}{3.4} there is exactly one such trivialization of $E'$ for
each given biextension $E$.  One can of course also prove directly
that (\SGARefLink{sga7.sga7i.eviii.n3-5-2}{3.5.2}) is an isomorphism for $i=0,1$ - which is equivalent to
the conclusion of \SGARefLink{sga7.sga7i.eviii.n3-5}{3.5} - by using the exact sequence of the
$\Biext^i$ in VII \SGARefLink{sga7.sga7i.evii.n3-7-5}{3.7.5} associated with
\[
 0\longrightarrow Q'\longrightarrow Q\longrightarrow Q''
 \longrightarrow0.
\]

\medskip
\noindent\underline{Remark \SGARefTarget{sga7.sga7i.eviii.n3-6}{3.6}.}  The symmetric statement obtained
from \SGARefLink{sga7.sga7i.eviii.n3-5}{3.5} by interchanging left and right is of course also valid; it
also follows formally from \SGARefLink{sga7.sga7i.eviii.n3-5}{3.5} by passing to symmetric biextensions
(VII \SGARefLink{sga7.sga7i.evii.n2-7}{2.7}).  Suppose, in addition to the data and hypotheses of \SGARefLink{sga7.sga7i.eviii.n3-5}{3.5},
that the symmetric data and hypotheses are given: namely, that $Q$
has connected fibres (equivalently, that $Q''$ has connected fibres,
since $Q'$ has connected fibres), and that there is an exact sequence
\[
 0\longrightarrow P'\longrightarrow P\longrightarrow P''
 \longrightarrow0
\]
with $P'$ a torus.  Successive application of \SGARefLink{sga7.sga7i.eviii.n3-5}{3.5} and of its
symmetric form, applied to this last sequence and to $Q''$, shows
that the functor
\begin{equation}\tag{\SGARefTarget{sga7.sga7i.eviii.n3-6-1}{3.6.1}}
 \operatorname{BIEXT}(P'',Q'';\underline G_m)
 \xrightarrow{\ \sim\ }
 \operatorname{BIEXT}(P,Q;\underline G_m)
\end{equation}
is an equivalence of categories.  In particular, the following
homomorphism is an isomorphism:
\begin{equation}\tag{\SGARefTarget{sga7.sga7i.eviii.n3-6-2}{3.6.2}}
 \Biext^i(P'',Q'';\underline G_m)
 \xrightarrow{\ \sim\ }
 \Biext^i(P,Q;\underline G_m),
 \qquad i=0,1.
\end{equation}

\medskip
\noindent\underline{Proposition \SGARefTarget{sga7.sga7i.eviii.n3-7}{3.7}.}  \underline{Let $A$ be an
abelian scheme over $S$, let $T$ be a group scheme of multiplicative
type (SGA 3 IX) and of finite type over $S$, and let}
\[
 \underline M=\uHom(T,\underline G_m)
\]
\underline{be the Cartier dual of $T$ (a finitely generated twisted
constant group over $S$, by \SGARefLink{sga7.sga7i.eviii.n3-3-3}{3.3.3}).  Then the categories
$\operatorname{EXT}(A,T)$ and
$\operatorname{BIEXT}(A,\underline M;\underline G_m)$, formed
respectively by commutative extensions of $A$ by $T$ and by
biextensions of $(A,\underline M)$ by $\underline G_m$, are rigid and
canonically equivalent.  Their groups of isomorphism classes are
canonically isomorphic to $\Hom(\underline M,A^*)$, where}
\[
 A^*=\uExt^1(A,\underline G_m)
\]
\underline{is the dual abelian scheme of $A$ (\SGARefLink{sga7.sga7i.eviii.n3-2}{3.2}).  In particular,
there is a canonical isomorphism}
\begin{equation}\tag{\SGARefTarget{sga7.sga7i.eviii.n3-7-1}{3.7.1}}
 \Ext^1(A,T)\xrightarrow{\ \sim\ }
 \Hom(\underline M,A^*).
\end{equation}

This is the special case of \SGARefLink{sga7.sga7i.eviii.n1-6}{1.6} obtained by taking $P=A$ and
$Q=\underline M$.  By the biduality theorem (SGA 3 X 5.7), we do
indeed have
\[
 T\simeq\uHom(\underline M,\underline G_m),
\]
and it remains only to verify the conditions (\SGARefLink{sga7.sga7i.eviii.n1-6-1}{1.6.1}), which here read
\[
 \uHom(A,\underline G_m)=0,
 \qquad
 \Hom\bigl(A,\uExt^1(\underline M,\underline G_m)\bigr)=0.
\]
The first equality was recalled in (\SGARefLink{sga7.sga7i.eviii.n3-2-1}{3.2.1}); the second may be
strengthened to
\[
 \uExt^1(\underline M,\underline G_m)=0.
\]
This assertion is local on $S$, so we may suppose that $\underline
M$ is constant, $\underline M=M_S$.  Decomposing $M$ as a direct sum
of cyclic $\mathbb Z$-modules reduces us to $M=\mathbb Z$, which is
trivial, and to $M=\mathbb Z/n\mathbb Z$, which follows from the
$n$-divisibility of $\underline G_m$ (or, more learnedly, from
\SGARefLink{sga7.sga7i.eviii.n3-3-1}{3.3.1}).

When $T=\underline G_m$, and hence $\underline M=\mathbb Z_S$,
(\SGARefLink{sga7.sga7i.eviii.n3-7-1}{3.7.1}) reduces to the identity isomorphism (\SGARefLink{sga7.sga7i.eviii.n3-2-3}{3.2.3}).  Of course, a
direct proof of (\SGARefLink{sga7.sga7i.eviii.n3-7-1}{3.7.1}) would be essentially trivial, without
introducing the group
\[
 \Ext^1\bigl(A\mathbin{\mathop{\otimes}^{\mathbb L}}\underline M,
              \underline G_m\bigr)
 =\Biext^1(A,\underline M;\underline G_m).
\]

